\section{Derivative Theorems and Methods}
\subsection{Differentiation Theorems and Methods}
\subsubsection{Linearity, or sum rule, difference rule, and constant multiple rule}
If functions $f$ and $g$ are differentiable on $I$ and $a,b\in\mathbb{R}$, then $af(x)+bg(x)$ is also differentiable on $I$ and its derivative is given by
\[\dv{}{x}\qty(af(x)+bg(x))=af'(x)+bg'(x).\]
\begin{proof}
\[\ba
\dv{}{x}\qty(af(x)+bg(x))&=\lim_{h\to 0}\frac{af(x+h)+bg(x+h)-af(x)-bg(x)}{h}\\
&=af'(x)+bg'(x)
\ea\]
\end{proof}
\subsubsection{Product rule (乘法定則)}
If functions $f$ and $g$ are differentiable on $I$, then the product $fg$ is also differentiable on $I$ and its derivative is given by
\[\dv{x}\qty(f(x)g(x))=f'(x)g(x)+f(x)g'(x).\]
\begin{proof}
\[\begin{aligned}
\dv{f(x)g(x)}{x}&=\lim_{h\to 0}\frac{f(x+h)g(x+h)-f(x)g(x)}{h}\\
&=\lim_{h\to 0}\frac{f(x+h)(g(x+h)-g(x))+g(x)(f(x+h)-f(x))}{h}\\
&=f(x)g'(x)+f'(x)g(x).
\end{aligned}\]
\end{proof}
\subsubsection{General Leibniz rule}
If functions $f$ and $g$ are $n$-times differentiable on $I$, then the product $fg$ is also $n$-times differentiable on $I$ and its derivative is given by
\[(fg)^{(n)}=\sum _{k=0}^{n}{n \choose k}f^{(n-k)}g^{(k)}.\]
\subsubsection{Product rule for dot product}
If functions $f$ and $g$ are differentiable on $I\subseteq\bbR^n$, then the dot product $f\cdot g$ is also differentiable on $I$ and its derivative is given by
\[\dv{x}\qty(f(x)\cdot g(x))=f'(x)\cdot g(x)+f(x)\cdot g'(x).\]
\begin{proof}
\[\begin{aligned}
\dv{f(x)\cdot g(x)}{x}&=\lim_{h\to 0}\frac{f(x+h)\cdot g(x+h)-f(x)\cdot g(x)}{h}\\
&=\lim_{h\to 0}\frac{f(x+h)\cdot (g(x+h)-g(x))+g(x)\cdot (f(x+h)-f(x))}{h}\\
&=f(x)\cdot g'(x)+f'(x)\cdot g(x).
\end{aligned}\]
\end{proof}
\subsubsection{Product rule for cross product}
If functions $f$ and $g$ are differentiable on $I\subseteq\bbR^3$, then the cross product $f\times g$ is also differentiable on $I$ and its derivative is given by
\[\dv{x}\qty(f(x)\times g(x))=f'(x)\times g(x)+f(x)\times g'(x).\]
\begin{proof}
\[\begin{aligned}
\dv{f(x)\times g(x)}{x}&=\lim_{h\to 0}\frac{f(x+h)\times g(x+h)-f(x)\times g(x)}{h}\\
&=\lim_{h\to 0}\frac{f(x+h)\times (g(x+h)-g(x))+g(x)\times (f(x+h)-f(x))}{h}\\
&=f(x)\times g'(x)+f'(x)\times g(x).
\end{aligned}\]
\end{proof}
\subsubsection{Quotient rule (除法定則)}
If functions $f$ and $g$ are differentiable on $I$ and $g\neq 0$ on $I$, then the quotient $\frac{f}{g}$ is also differentiable on $I$ and its derivative is given by
\[\dv{x}\qty(\frac{f(x)}{g(x)})=\frac{f'(x)g(x)-f(x)g'(x)}{\qty(g(x))^2}.\]
\begin{proof}
\[\begin{aligned}
\dv{}{x}\frac{f(x)}{g(x)}&=\lim_{h\to 0}\frac{\frac{f(x+h)}{g(x+h)}-\frac{f(x)}{g(x)}}{h}\\
&=\lim_{h\to 0}\frac{f(x+h)g(x)-f(x)g(x+h)}{h\,g(x+h)g(x)}\\
&=\lim_{h\to 0}\frac{(f(x+h)-f(x))g(x)-f(x)(g(x+h)-g(x))}{h\,g(x+h)g(x)}\\
&=\lim_{h\to 0}\frac{f(x+h)-f(x)}{h\,g(x+h)}-f(x)\lim_{h\to 0}\frac{g(x+h)-g(x)}{h\,g(x+h)g(x)}\\
&=\frac{f'(x)g(x)-f(x)g'(x)}{g(x)^2}
\end{aligned}\]
\end{proof}
\subsubsection{Chain rule (連鎖律) of real functions}
Let $f\colon X\subseteq\bbR\to\bbR$ and $g\colon Y\subseteq\bbR\to\bbR$ be functions, $g$ be differentiable on $I\subseteq\bbR$, and $f$ be differentiable on $g(I)$. Then the composite $f\circ g$ is differentiable on $I$ and its derivative on $I$ is given by
\[\dv{}{x}\qty((f\circ g)(x))=f'\qty(g(x))g'(x).\]
\begin{proof}
\[\ba
\dv{}{x}\qty((f\circ g)(x))&=\lim_{h\to 0}\frac{f\qty(g(x+h))-f\qty(g(x))}{h}\\
&=\lim_{h\to 0}\frac{f\qty(g(x+h))-f\qty(g(x))}{g(x+h)-g(x)}\cdot\frac{g(x+h)-g(x)}{h}\\
&=f'\qty(g(x))g'(x)
\ea\]
\end{proof}
\sssc{Multivariable chain rule}
Suppose
\[F\qty(t,\mb{s}_1,\mb{s}_2,\ldots,\mb{s}_n)=f\qty(x_1(t,\mb{s}_1),x_2(t,\mb{s}_2),\ldots,x_n(t,\mb{s}_n)),\]
then
\[\pdv{F}{t}=\sum_{i=1}^n\pdv{f}{x_i}\pdv{x_i}{t}.\]
\subsubsection{Faà di Bruno's formula}
\bma
&\frac{\mathrm{d}^n}{\mathrm{d}x}f\left(g(x)\right)\\
=&\sum_{\sum_{i=1}^nim_i=n,\quad m_i\in\mathbb{N}_0}\frac{n!}{\prod_{j=1}^nm_j!}\cdot\\
&f^{\qty(\sum_{j=1}^nm_j)}\left(g(x)\right)\cdot\\
&\prod_{j=1}^n\left(\frac {g^{(j)}(x)}{j!}\right)^{m_j}.
\eam
\subsubsection{Chain rule of real vector functions}
Let $f\colon X\subseteq\bbR^n\to\bbR^p$ and $g\colon Y\subseteq\bbR^m\to\bbR^n$ be functions, $g$ be differentiable on $\Omega\subseteq\bbR^m$, and $f$ be differentiable on $g(\Omega)$. Then the composite $f\circ g$ is differentiable on $\Omega$ and its derivative on $\Omega$ is given by
\[\nabla\qty(f\circ g)(\mb{x})=(\nabla f)(g(\mb{x}))\cdot(\nabla g)(\mb{x}).\]
\subsubsection{Exponential Bell polynomials}
Partially or incomplete exponential Bell polynomials $B_{n,k}\qty(x_1,x_2,\ldots, x_{n-k+1})$ is defined as
\[B_{n,k}\qty(x_1,x_2,\ldots, x_{n-k+1})=n!\sum_{\langle j_i\rangle\in C}\prod_{i=1}^{n-k+1}\frac{x_i^{\phantom{i}j_i}}{(i!)^{j_i}j_i!},\]
where
\[C=\left\{\langle j_i\rangle_{i=1}^{n-k+1}\middle|\sum_{i=1}^{n-k+1}=k\land\sum_{i=1}^{n-k+1}ij_i=n\right\}.\]
Complete exponential Bell polynomial $B_n\qty(x_1,x_2,\ldots, x_n)$ is defined as
\[B_n\qty(x_1,x_2,\ldots, x_n)=n!\sum_{k=0}^nB_{n,k}\qty(x_1,x_2,\ldots, x_{n-k+1}).\]
\subsubsection{Derivative of inverse functions}
\[\qty(f^{-1})'(x)=\frac{1}{f'\qty(f^{-1}(x)).\]
\begin{proof}
\[f\qty(f^{-1}(x))=x.\]
Differentiate both sides:
\[f'\qty(f^{-1}(x))\qty(f^{-1})'(x)=1.\]
\end{proof}
\[\qty(f^{-1})''(x)=\frac{f''\qty(f^{-1}(x))}{\qty(f'\qty(f^{-1}(x)))^3}.\]
\begin{proof}
\[f'\qty(f^{-1}(x))\qty(f^{-1})'(x)=1.\]
Differentiate both sides:
\[f''\qty(f^{-1}(x))\qty(\qty(f^{-1})'(x))^2+f'\qty(f^{-1}(x))\qty(f^{-1})''(x)=0.\]
\[\qty(f^{-1})''(x)=\frac{f''\qty(f^{-1}(x))\qty(\qty(f^{-1})'(x))^2}{f'\qty(f^{-1}(x))}=\frac{f''\qty(f^{-1}(x))}{\qty(f'\qty(f^{-1}(x)))^3}.\]
\end{proof}
\[\dv[n]{f^{-1}(x)}{x}=\frac{(-1)^{n-1}}{\qty(f'\qty(f^{-1}(x)))^{2n-1}}B_{n-1}\sum_{i=2}^{n}f^{(i)}\qty(f^{-1}(x)),\quad n\in\bbN,\]
where $B_n$ is complete exponential Bell polynomial.
\sssc{Jacobian of inverse}
Let $\mathbf{T}\colon I\subseteq\mathbb{R}^n\to\mathbb{R}^n$ be injective and differentiable on $D\subseteq I$, and $\mb{T}^{-1}$ be differentiable on $\mb{T}(D)$. Then,
\[J(\mb{T})J(\mb{T}^{-1})=\mb{I}\]
where $\mb{I}$ is the $n\times n$ unit matrix.
\bpr
PLACEHOLDER
\epr
\subsubsection{Implicit function theorem (隱函數定理) of two real variables}
Given a function $f(x,y)$ is a function that is continuously differentiable on a neighborhood of a point $(x_0,y_0)$, $f(x_0,y_0)=0$, and $\pdv{f}{y}(x_0,y_0)\neq 0$ on a neighborhood of $(x_0,y_0)$, then there exists a unique differentiable function $g$ such that $y_0=g(x_0)$ and $f(x,g(x))=0$ on a neighborhood of $x_0$. Moreover, by differentiating both sides of $f(x,g(x))=0$, called implicit differentiation, we obtain
\[g'(x)=-\frac{\pdv{f}{x}(x,g(x))}{\pdv{f}{y}(x,\varphi(x))}.\]
\subsubsection{Implicit function theorem of multiple real variables}
Below, we will denote $(\mb{x},\mb{y})$ as an element of $\bbR^{(n+m)}$ for some $\mb{x}\in\bbR^n$ and $\mb{y}\in\bbR^m$, and denote the left $n$ columns of $(Df)(\mb{a},\mb{b})$ as $J_{f,\mb{x})$, and the right $m$ columns of $(Df)(\mb{a},\mb{b})$ as $J_{f,\mb{y})$. Given a function $f\colon D\subseteq\bbR^{(n+m)}\to\bbR^m$ that is continuously differentiable on a neighborhood of a point $(\mb{a},\mb{b})$, $f(\mb{a},\mb{b})=\mb{0}$, and $J_{f,\mb{y})(\mb{a},\mb{b})$ is invertible, then there exists a unique function $g$ such that $\mb{b}=g(\mb{a})$ and $f(\mb{x},g(\mb{x}))=\mb{0}$ on a neighborhood of $\mb{a}$. Moreover, by differentiating both sides of $f(\mb{x},g(\mb{x}))=\mb{0}$, called implicit differentiation, we obtain
\[(Dg)(\mb{x})=-\qty[J_{f,\mb{y}}(\mb{x},g(\mb{x}))]^{-1}\qty[J_{f,\mb{x}}(\mb{x},g(\mb{x}))],\]
or for each component $x_i$ of $\mb{x}$,
\[\pdv{g}{x_i}=-\qty[J_{f,\mb{y}}(\mb{x},g(\mb{x}))]^{-1}\pdv{f}{x_i}.\]
\subsubsection{Derivative of natural logarithm of absolute value of function}
If functions $f$ is differentiable and nonzero on $I$, then $\ln|f(x)|$ is also differentiable on $I$ and its derivative is given by
\[\dv{\ln|f(x)|}{x}=\frac{f'(x)}{f(x)}.\]
\subsubsection{Logarithmic differentiation}
Logarithmic differentiation is a method for finding derivatives by taking the natural logarithm of (the absolute values of) both sides of an equation first, and then differentiating.

It’s especially useful when the function involves products or quotients of many factors, or variables in both the base and the exponent.
\subsubsection{Derivative of variable base, constant exponent funcion}
If $n\in\bbR$ and $f(x)$ is differentiable on $I$, then $\qty(f(x))^n$ is also differentiable on $I$ and its derivative is given by
\[\dv{}{x}\qty(f(x))^n=n\qty(f(x))^{n-1}f'(x).\]
\subsubsection{Derivative of constant base, variable exponent funcion}
If $n\in\bbR$ and $f(x)$ is differentiable on $I$, then $n^{f(x)}$ is also differentiable on $I$ and its derivative is given by
\[\dv{}{x}n^{f(x)}=n^{f(x)}\qty(\ln n)f'(x).\]
\subsubsection{Derivative of variable base, variable exponent funcion}
\[\dv{}{x}\qty(f(x))^{g(x)}=\qty(f(x))^{g(x)}\qty(g'(x)\ln(f(x))+g(x)\frac{f'(x)}{f(x)}).\]
\begin{proof}
\[y\coloneq\qty(f(x))^{g(x)}\]
Take natural logarithm of both sides:
\[\ln y=g(x)\ln(f(x)).\]
Take differentiation of both sides:
\[\frac{y'}{y}=g'(x)\ln(f(x))+g(x)\frac{f'(x)}{f(x)}.\]
\[y'=\qty(f(x))^{g(x)}\qty(g'(x)\ln(f(x))+g(x)\frac{f'(x)}{f(x)}).\]
\end{proof}
\sssc{Existence of total derivative}
For a function $f\colon D\subseteq\bbR^n\to\bbR^m$, $f'$ exists at a point $a\in D$ iff
\[\lim_{h\to 0}\frac{f(a+h)-L(a+h)}{\|h\|}=0,\]
where $L$ is the linearization of $f$ at $a$, that is,
\[L(x)=f(a)+\nabla f(a)(x-a),\]
and the derivative of $f$ at $a$ is $\nabla f(a)$.
\bpr
$f'(a)$ exists iff there exists a bounded linear operator $A\colon\bbR^n\to\bbR^m$ such that
\[\lim_{h\to0}\frac{f(a+h)-f(a)-A(h)}{\|h\|}=0.\]
Suppose the limit is 0. For each $x_i$, let the unit vector along its positive direction be $e_i$, along the line $x_it,\quad t\in\bbR$,
\[\lim_{t\to0}\frac{f(a+e_it)-f(a)-A(e_it)}{t}=\pdv{f}{x_i}(a)-A(e_i)\]
Thus $A_i=\pdv{f}{x_i}(a)$, and
\[A=\nabla f(a)\]
And
\[\frac{f(a+h)-f(a)-A(h)}{\|h\|}=\frac{f(a+h)-f(a)+\nabla f(a)(h)}{\|h\|}=\frac{f(a+h)-L(a+h)}{\|h\|}\]
\epr
\subsection{Graph-related}
\subsubsection{Differentiability implies continuity}
Let $V$ and $W$ be normed vector spaces, $U$ be an open subset of $V$, and $f\colon U\subseteq V\to W$ be a function. For a point $a\in U$, $f$ is differentiable at $a$ implies it is continuous at $a$; for an open subset $I\subseteq U$, $f$ is differentiable on $I$ implies it is continuous on $I$.
\begin{proof}
$f$ is differentiable at $a$ iff there exists bounded linear operator $A\colon V\to W$ such that
\[\lim_{\|h\|_V\to 0}\frac{\|f(a+h)-f(a)-A(h)\|_W}{\|h\|_V}=0,\]
that is,
\[\forall\varepsilon>0\colon\exists\delta>0\text{\ s.t.\ }\|h\|_V<\delta\implies\|f(a+h)-f(a)-A(h)\|_W\leq\varepsion\|h\|_V.\]
By the triangular identity,
\[\|f(a+h)-f(a)\|_W\leq\|A(h)\|_W+\|f(a+h)-f(a)-A(h)\|_W.\]
By the definition of bounded linear operator, there exists $C\in\bbR_{>0}$ such that
\[\|A(h)\|_W\leq C\|h\|_V.\]
\[\|f(a+h)-f(a)\|_W\leq (C+\varepsilon)\|h\|_V.\]
Thus,
\[\forall\varepsilon>0\colon\exists\delta'=\frac{\varepsilon}{C+\varepsilon}\text{\ s.t.\ }\|h\|_V<\delta'\implies\|f(a+h)-f(a)\|_W\leq(C+\varepsilon)\|h\|_V<(C+\varepsilon)\delta'=\varepsilon.\]
\end{proof}
\subsubsection{Extreme Value Theorem (EVT) (極值定理)}
Let function $f\colon I\subseteq\mathbb{R}\to\mathbb{R}$ be continuous on an interval $[a,b]\subseteq I$, then
\[\exists c,d \in [a, b] \text{\ s.t.\ }\forall x \in [a, b]\colon f(c)\geq f(x)\geq f(d).\]
\begin{proof}
We first claim that $f$ is bounded on $[a,b]$. Argue by contradiction. Suppose $f$ is not bounded-above.

For any $n\in\mathbb{N}$, there exists $x_n\in [a,b]$ such that $f\qty(x_n)>n$. This defines a sequence $\langle x_n\rangle$.

By the Bolzano-Weierstrass Theorem and that $[a,b]$ is bounded, it has a convergent subsequence $x_{n_k}$ that converges to some limit $L$. Since $[a,b]$ is closed, $L\in [a,b]$.

Since $f$ is continuous on $[a,b]$, the sequence $f\qty(x_{n_k})$ converges to $f(L)$, which is finite.

This contradicts $f\qty(x_n)>n$ for every $n\in\mathbb{N}$. So the set
\[S=\{f(x)\mid x\in [a,b]\}.\]
must be bounded-above.

By considering $-f$, we obtain that $S$ is also bounded-below.

By the Completeness of the Real Numbers, every non-empty set bounded-above has a least upper bound; every non-empty set bounded-below has a greatest lower bound.

Let
\[M=\sup S,\quad m=\inf S.\]

Since $M$ is the least upper bound of $S$, for any $n\in\mathbb{N}$, $M-\frac{1}{n}$ is not an upper bound of $S$. This means that for each $n\in\mathbb{N}$, there exists a point $y_n\in [a,b]$ such that:
\[M-\frac{1}{n}<f(y_n)\leq M.\]
This defines another sequence $\langle y_n\rangle$ within $[a, b]$.

By the Bolzano-Weierstrass Theorem and that $[a,b]$ is bounded, it has a convergent subsequence $y_{n_k}$ that converges to some limit $c$. Since $[a,b]$ is closed, $c\in [a,b]$.
\[\lim_{n_k\to \infty}\qty(M - \frac{1}{n_k}) = M.\]
By the Squeeze Theorem:
\[\lim_{n_k \to \infty} f\qty(y_{n_k}) = M.\]
Since $f$ is continuous on $[a,b]$, the sequence $f\qty(y_{n_k})$ converges to $f(c)$.

Similarly, we can prove that there exists $d\in [a,b]$ such that $f(d)=m$.
\end{proof}
\subsubsection{Extreme Value Theorem (EVT) for compact set}
Let $X$ be a topological space, and function $f\colon I\subset X\to\mathbb{R}$ be continuous on a compact set $D\subseteq I$, then
\[\exists a,b \in D \text{\ s.t.\ }\forall x \in D\colon f(a)\geq f(x)\geq f(b).\]
\subsubsection{Fermat's Theorem (費馬引理) or Interior Extremum Theorem (內部極值定理)}
Let $X$ be a LCTVS and function $f\colon D\subseteq X\to\mathbb{R}$ be differentiable on an open set $U\subseteq D$ and $c\in U$ be a local extreme of $f$, then
\[f'(c)=0,\]
where $f'(c)\in L(X,\bbR)$ and $0$ is the zero element of $L(X,\bbR)$.
\begin{proof}
Let $h\in X$ be arbitrary.

Because $U$ is open, there exists $\varepsilon>0$ such that $c+th\in U$ for $\abs{t}<\varepsilon$. Define
\[g(t)=f(c+th),\quad\abs{t}<\varepsilon.\]
Since $f$ is differentiable on $U$, $g$ is differentiable and
\[g'(0)=f'(c)(h).\]
Without loss of generality, suppose $c$ is a local maximum. Then, there exists $\delta>0$ such that
\[g(t)\leq g(0)\]
for all $\abs{t}<\delta$.

Thus for all $\delta>t>0$,
\[\frac{g(t)-g(0)}{t}\leq 0\]
and for all $-\delta<t<0$,
\[\frac{g(t)-g(0)}{t}\geq 0.\]
Therefore
\[g'(0)=0.\]
Since
\[g'(0)=f'(c)(h)=0\]
for all $h\in X$
\[f'(c)=0.\]
\end{proof}
In terms of critical point, Fermat's Theorem can be rephrased as, a local extreme is a critical point.
\subsubsection{Rolle's Theorem (羅爾定理)}
Let function $f\colon I\subseteq\mathbb{R}\to\mathbb{R}$ be continuous on an interval $[a, b]\subseteq I$ and differentiable on $(a,b)$, and $f(a)=f(b)$, then
\[\exists c\in (a, b)\text{\ s.t.\ }f'(c)=0.\]
\begin{proof}
By the Extreme Value Theorem, $f$ must have maximum and minimum on $[a,b]$.

Let
\[M=\max_{x\in [a,b]}f(x),\quad m=\min_{x\in [a,b]}f(x).\]
If $M=m$, then $f$ is constant on $[a,b]$, then $f'(x)=0$ for any $c\in (a,b)$.

Otherwise, $M>m$. Because $f(a)=f(b)$. 

We claim that there must be $c\in (a,b)$ such that $M=f(c)$ or $m=f(c)$. Argue by contradiction. If the claim is false, then $M=m=f(a)=f(b)$, contradicting $M>m$.

By Fermat's Theorem,
\[f'(c)=0.\]
\end{proof}
\subsubsection{Mean Value Theorem (MVT) (均值定理) (for Derivatives)}
Let function $f\colon I\subseteq\mathbb{R}\to\mathbb{R}$ be continuous on an interval $[a, b]\subseteq I$ and differentiable on $(a,b)$ with $a<b$, then
\[\exists c\in (a, b)\text{\ s.t.\ }f'(c)=\frac{f(b)-f(a)}{b-a}.\]
\begin{proof}
Let
\[\phi(x) = f(x) - \frac{f(b)-f(a)}{b-a} (x-a).\]
\[\phi(a)=\phi(b)=f(a).\]
By Rolle's Theorem, there exists $c\in (a,b)$ such that $\phi'(c)=0$.
\[\phi'(c)=f'(c)-\frac{f(b)-f(a)}{b-a}.\]
\end{proof}
\subsubsection{Cauchy's Mean Value Theorem (CMVT) (柯西均值定理)}
Let functions $f\colon I\subseteq\mathbb{R}\to\mathbb{R}$ and $g\colon J\subseteq\mathbb{R}\to\mathbb{R}$ be continuous on an interval $[a, b]\subseteq I\cap J$ and differentiable on $(a,b)$ with $a<b$, then
\[\exists c\in (a,b)\text{\ s.t.\ }\qty(f(b)-f(a))g'(c)=\qty(g(b)-g(a))f'(c).\]
\begin{proof}
    If $f(b)-f(a)=g(b)-g(a)=0$, it holds. For the other cases, without loss of generality, assume $g(b)-g(a)\neq 0$.

    Define a function $\phi(x)$
    \[\phi(x)=f(x)-\frac{f(b)-f(a)}{g(b)-g(a)}g(x).\]
    \[\phi(b)-\phi(a)=\qty(f(b)-f(a))-\frac{f(b)-f(a)}{g(b)-g(a)}\qty(g(b)-g(a))=0.\]
    By Rolle's Theorem, there exists $c\in (a,b)$ such that $\phi'(c)=0$.
    \[\phi'(x)=f'(x)-\frac{f(b)-f(a)}{g(b)-g(a)}g'(x).\]
\end{proof}
\subsubsection{Generalized Racetrack principle}
If $f'(x)>g'(x)$ for all $x\in (a,b)$ with $b\in\bbR_{>a}\cup\{\infty\}$, and $f(a)=g(a)$, then $f(x)>g(x)$ for all $x\in(a,b)$.

If $f'(x)>g'(x)$ for all $x\in (a,b)$ with $b\in\bbR_{>a}\cup\{\infty\}$, $f$ is continuous on $(a,b]$, and $f(a)=g(a)$, then $f(x)>g(x)$ for all $x\in(a,b]$.

If $f'(x)\geq g'(x)$ for all $x\in (a,b)$ with $b\in\bbR_{>a}\cup\{\infty\}$, and $f(a)=g(a)$, then $f(x)\geq g(x)$ for all $x\in(a,b]$.

If $f'(x)\geq g'(x)$ for all $x\in (a,b)$ with $b\in\bbR_{>a}\cup\{\infty\}$, $f$ is continuous on $(a,b]$, and $f(a)=g(a)$, then $f(x)\geq g(x)$ for all $x\in(a,b]$.
\begin{proof}
    For the first and second statements, define a function
    \[h(x)=f(x)-g(x).\]
    \[h'(x)=f'(x)-g'(x)>0,\quad\forall x\in(a,b).\]
    \[h(a)=f(a)-g(a)=0.\]
    By MVT, for all $x\in(a,b)$ for the first statement and $x\in(a,b]$ for the second statement, there exists $c\in(a,x)$ such that
    \[\frac{h(x)-h(a)}{x-a}=h'(c)>0.\]
    \[f(x)-g(x)=h(x)=(x-a)h'(c)>0.\]
    The other statements can be proven similarly.
\end{proof}
\subsubsection{Increasing/Decreasing Test (I/D Test) Theorem}
Let function $f\colon I\subseteq\bbR\to\bbR$ be differentiable on an open interval $(a,b)\subseteq I$.
\bit
\item if $f'(x)>0$ for any $x\in (a,b)$, then $f$ is strictly increasing on $(a,b)$.
\item If $f'(x)>0$ for any $x\in (a,b)$ and $f$ is continuous on $[a,b)$, then $f$ is strictly increasing on $[a,b)$.
\item If $f'(x)>0$ for any $x\in (a,b)$ and $f$ is continuous on $(a,b]$, then $f$ is strictly increasing on $(a,b]$.
\item if $f'(x)<0$ for any $x\in (a,b)$, then $f$ is strictly decreasing on $(a,b)$.
\item If $f'(x)<0$ for any $x\in (a,b)$ and $f$ is continuous on $[a,b)$, then $f$ is strictly decreasing on $[a,b)$.
\item If $f'(x)<0$ for any $x\in (a,b)$ and $f$ is continuous on $(a,b]$, then $f$ is strictly decreasing on $(a,b]$.
\item $f'(x)\geq 0$ for any $x\in (a,b)$ iff $f$ is non-decreasing on $(a,b)$.
\item If $f'(x)\geq 0$ for any $x\in (a,b)$ and $f$ is continuous on $[a,b)$, then $f$ is non-decreasing on $[a,b)$.
\item If $f'(x)\geq 0$ for any $x\in (a,b)$ and $f$ is continuous on $(a,b]$, then $f$ is non-decreasing on $(a,b]$.
\item $f'(x)\leq 0$ for any $x\in (a,b)$ iff $f$ is non-increasing on $(a,b)$.
\item If $f'(x)\leq 0$ for any $x\in (a,b)$ and $f$ is continuous on $[a,b)$, then $f$ is non-increasing on $[a,b)$.
\item If $f'(x)\leq 0$ for any $x\in (a,b)$ and $f$ is continuous on $(a,b]$, then $f$ is non-increasing on $(a,b]$.
\eit
\begin{proof}
    For the first statement, we want to show that $\forall c,d\in (a,b)\colon c<d\implies f(c)<f(d)$ if $f'(x)>0$ for any $x\in (a,b)$.

    Take any two points $c,d\in (a,b)$ with $c<d$. By MVT, there exists $g\in (c,d)$ such that
    \[f'(g)=\frac{f(d)-f(c)}{d-c}.\]
    Because $f'(g)>0$ and $d-c>0$, we get $f(d)-f(c)>0$.
    
    The second and third statements are obviously true as the first holds.
    
    The "if" directions of the other statements can be proven similarly.
    
    For the "only if" direction of the seventh statement, we want to show that $f'(x)\geq 0$ if $\forall c,d\in (a,b)\colon c<d\implies f(c)\leq f(d)$ and $f$ is differentiable on $(a,b)$.
    
    Fix any $x\in (a,b)$. For any $h>0$ with $(x+h)\in(a,b)$, we have
    \[f(x)\leq f(x+h).\]
    For any $h<0$ with $(x+h)\in(a,b)$, we have
    \[f(x)\geq f(x+h).\]
    Thus for any $h$ with $(x+h)\in(a,b)$, we have
    \[\frac{x+h}-f(x)}{h}\geq 0.\]
    The "only if" direction of the tenth statement can be proven similarly.
\end{proof}
\subsubsection{First Derivative Test for Local Extrema Theorem}
Let $f\colon I\subseteq\mathbb{R}\to\mathbb{R}$ be a function, $c$ be a critical point of $f$, and $f'$ be continuous at $c$.
\bit
\item If exist $a<c<b$ such that $f'(c)>0$ for any $x\in (a,c)$ and $f'(c)<0$ for any $x\in (c,b)$, then $f$ has a strict local maximum at $c$.
\item If exist $a<c<b$ such that $f'(c)\geq 0$ for any $x\in (a,c)$ and $f'(c)\leq 0$ for any $x\in (c,b)$, then $f$ has a local maximum at $c$.
\item If exist $a<c<b$ such that $f'(c)<0$ for any $x\in (a,c)$ and $f'(c)>0$ for any $x\in (c,b)$, then $f$ has a strict local minimum at $c$.
\item If exist $a<c<b$ such that $f'(c)\leq 0$ for any $x\in (a,c)$ and $f'(c)\geq 0$ for any $x\in (c,b)$, then $f$ has a local minimum at $c$.
\item If exist $a<c<b$ such that $f'(c)>0$ for any $x\in (a,b)$, then $f$ has no local exteme at $c$.
\item If exist $a<c<b$ such that $f'(c)<0$ for any $x\in (a,b)$, then $f$ has no local ext
eme at $c$.
\eit
\begin{proof}
    By Increasing/Decreasing Test Theorem, it is obvious.
\end{proof}
\subsubsection{Second derivative test of one variable}
Let $f\colon I\subseteq\mathbb{R}\to\mathbb{R}$ be a function, $c$ be a critical point of $f$, and $f$ be twice differentiable at $c$.
\bit
\item If $f''<0$, then $f$ has a strict local maximum at $c$.
\item If $f''>0$, then $f$ has a strict local minimum at $c$.
\eit
\subsubsection{Second (partial) derivative test of two variables}
Let $f(x,y)\colon I\subseteq\mathbb{R}^2\to\mathbb{R}$ be a function, $c$ be a critical point of $f$, $f$ be twice differentiable at $c$, and
\[D=\begin{vmatrix}
f_{xx} & f_{xy}\\
f_{yx} & f_{yy}
\end{vmatrix}.\]
\bit
\item If $D(c)>0$ and $f_{xx}(c)<0$, then $f$ has a strict local maximum at $c$.
\item If $D(c)>0$ and $f_{xx}(c)>0$, then $f$ has a strict local minimum at $c$.
\item If $D(c)<0$, then $f$ has a saddle point at $c$.
\eit
\subsubsection{Second derivative test of multiple derivatives}
Let $f\colon I\subseteq\mathbb{R}^n\to\mathbb{R}$ be a function, $c$ be a critical point of $f$, and $f$ be twice differentiable at $c$.
\bit
\item If $H(f)$ is negative definite, then $f$ has a strict local maximum at $c$.
\item If $H(f)$ is postive definite, then $f$ has a strict local minimum at $c$.
\item If $H(f)$ is negative semidefinite, then $f$ has a local maximum at $c$.
\item If $H(f)$ is postive semidefinite, then $f$ has a local minimum at $c$.
\item If $H(f)$ is indefinite, then $f$ has a saddle point at $c$.
\eit
\subsubsection{First Derivative Test for Concavity Theorem}
Let $f\colon I\subseteq\bbR\to\bbR$ be a function and $(a,b)\subseteq I$ be an open interval.
\bit
\item If $f'$ is strictly increasing on $(a,b)$, often called "$f$ lies strictly below all its tagent lines on $(a,b)$", then $f$ is strictly concave upward on $(a,b)$.
\item If $f'$ is strictly increasing on $(a,b)$ and $f$ is continuous on $[a,b)$, then $f$ is strictly concave upward on $[a,b)$.
\item If $f'$ is strictly increasing on $(a,b)$ and $f$ is continuous on $(a,b]$, then $f$ is strictly concave upward on $(a,b]$.
\item If $f'$ is strictly decreasing on $I$, often called "$f$ lies strictly above all its tagent lines on $(a,b)$", then $f$ is strictly concave downward on $(a,b)$.
\item If $f'$ is strictly decreasing on $(a,b)$ and $f$ is continuous on $[a,b)$, then $f$ is strictly concave downward on $[a,b)$.
\item If $f'$ is strictly decreasing on $(a,b)$ and $f$ is continuous on $(a,b]$, then $f$ is strictly concave downward on $(a,b]$.
\item $f'$ is non-decreasing on $(a,b)$, often called "$f$ lies weakly below all its tagent lines on $(a,b)$", iff $f$ is weakly concave upward on $(a,b)$.
\item If $f'$ is non-decreasing on $(a,b)$ and $f$ is continuous on $[a,b)$, then $f$ is weakly concave upward on $[a,b)$.
\item If $f'$ is non-decreasing on $(a,b)$ and $f$ is continuous on $(a,b]$, then $f$ is weakly concave upward on $(a,b]$.
\item $f'$ is non-increasing on $(a,b)$, often called "$f$ lies weakly above all its tagent lines on $(a,b)$", iff $f$ is weakly concave downward on $(a,b)$.
\item If $f'$ is non-increasing on $(a,b)$ and $f$ is continuous on $[a,b)$, then $f$ is weakly concave downward on $[a,b)$.
\item If $f'$ is non-increasing on $(a,b)$ and $f$ is continuous on $(a,b]$, then $f$ is weakly concave downward on $(a,b]$.
\eit
\begin{proof}
    For the first statement, we want to show that for any $c,d\in (a,b)$ and $\alpha\in(0,1)$, $f((1-\alpha)c+\alpha d)<(1-\alpha)f(c)+\alpha f(d)$ if for any $a\leq A<B\leq b$, $f'(A)<f'(B)$.

    Fix any $a\leq c<x<d\leq b$. By MVT, there exists $A\in (c,x)$ such that
    \[f'(A)=\frac{f(x)-f(c)}{x-c}.\]
    By MVT, there exists $B\in (b,c)$ such that
    \[f'(B)=\frac{f(d)-f(x)}{d-x}.\]
    Since $A<B$, $f'(A)\leq f'(B)$.
    \[\frac{f(x)-f(c)}{x-c}\leq\frac{f(d)-f(x)}{d-x}.\]
    \[(f(x)-f(c))(d-x)\leq(f(d)-f(x))(x-c).\]
    \[(d-c)f(x)\leq(d-x)f(c)+(x-c)f(d).\]
    \[f(x)\leq\frac{(d-x)}{(d-c)}f(c)+\frac{(x-c)}{(d-c)}f(d).\]
    Set
    \[\alpha=\frac{(x-c)}{(d-c)}\in(0,1).\]
    Then
    \[x=(1-\alpha)c+\alpha d.\]
    Thus the inequality becomes
    \[f\qty((1-\alpha)c+\alpha d)\leq(1-\alpha)f(c)+\alpha f(d).\]
    
    For the second statement, argue by contradiction. Assume that there exists $d\in(a,b)$ and $\alpha\in(0,1)$ such that
    \[f((1-\alpha)a+\alpha d)\geq (1-\alpha)f(a)+\alpha f(d).\]
    Let $x=(1-\alpha)a+\alpha d$.
    \[f(x)-f(a)\geq\alpha(f(d)-f(a)).\]
    By MVT, there exists $B\in (a,x)$ such that
    \[f(x)-f(a)=f'(B)(x-a)=f'(B)\alpha (d-a).\]
    By MVT, there exists $c\in (x,d)$ such that
    \[f(d)-f(x)=f'(C)(d-x)=f'(C)(1-\alpha)(d-a).\]
    \[f(d)-f(a)=(f'(B)\alpha+f'(C)(1-\alpha))(d-a).\]
    \[f(x)-f(a)=f'(B)\alpha (d-a)\geq\alpha(f'(B)\alpha+f'(C)(1-\alpha))(d-a).\]
    \[f'(B)\geq f'(B)\alpha+f'(C)(1-\alpha).\]
    \[f'(B)\geq f'(C).\]
    This contradicts the condition that $B<C\implies f'(B)<f'(C)$ for any $a<B<C<b$.
    
    The third statement can be proven similarly.
    
    The "if" directions of the other statements can be proven similarly.
    
    For the "only if" direction of the seventh statement, we want to show that for any $c,d\in (a,b)$ and $\alpha\in(0,1)$, $f((1-\alpha)c+\alpha d)\leq (1-\alpha)f(c)+\alpha f(d)$ if $\forall c,d\in (a,b)\colon c<d\implies f'(c)\leq f'(d)$.
    
    Fix any $c,d\in(a,b)$ with $c<d$.
    
    By MVT, there exists $g\in(c,((1-\alpha)c+\alpha d))$ such that
    \[f((1-\alpha)c+\alpha d)-f(c)=f'(g)((1-\alpha)c+\alpha d-c)=f'(g)\alpha(d-c).\]
    By MVT, there exists $h\in(((1-\alpha)c+\alpha d)),d)$ such that
    \[f(d)-f((1-\alpha)c+\alpha d)=f'(h)(d-(1-\alpha)c-\alpha d)=f'(h)(1-\alpha)(d-c).\]
    Thus,
    \[f((1-\alpha)c+\alpha d)=f(c)+f'(g)\alpha(d-c)=f(d)-f'(h)(1-\alpha)(d-c).\]
    \[f'(g)\alpha+f'(h)(1-\alpha)=\frac{f(d)-f(c)}{d-c}.\]
    Because $f'(g)\leq f'(h)$, we have
    \[f'(g)(d-c)\leq f(d)-f(c).\]
    Substitute:
    \[f((1-\alpha)c+\alpha d)\leq f(c)+\alpha(f(d)-f(c))=(1-\alpha)f(c)+\alpha f(d).\]
    The "only if" direction of the tenth statement can be proven similarly.
\end{proof}
\subsubsection{Concavity Test Theorem or Second Derivative Test for Concavity Theorem}
Let function $f\colon I\subseteq\bbR\to\bbR$ be twice differentiable on an open interval $(a,b)\subseteq I$.
\bit
\item If $f''>0$ for any $x\in (a,b)$, then the graph of $f$ is strictly concave upward on $(a,b)$.
\item If $f''>0$ for any $x\in (a,b)$ and $f'$ is continuous on $[a,b)$, then the graph of $f$ is strictly concave upward on $[a,b)$.
\item If $f''>0$ for any $x\in (a,b)$ and $f'$ is continuous on $(a,b]$, then the graph of $f$ is strictly concave upward on $(a,b]$.
\item If $f''<0$ for any $x\in (a,b)$, then the graph of $f$ is strictly concave downward on $(a,b)$.
\item If $f''<0$ for any $x\in (a,b)$ and $f'$ is continuous on $[a,b)$, then the graph of $f$ is strictly concave downward on $[a,b)$.
\item If $f''<0$ for any $x\in (a,b)$ and $f'$ is continuous on $(a,b]$, then the graph of $f$ is strictly concave downward on $(a,b]$.
\item $f''\geq 0$ for any $x\in (a,b)$ iff the graph of $f$ is weakly concave upward on $(a,b)$.
\item If $f''\geq 0$ for any $x\in (a,b)$ and $f'$ is continuous on $[a,b)$, then the graph of $f$ is weakly concave upward on $[a,b)$.
\item If $f''\geq 0$ for any $x\in (a,b)$ and $f'$ is continuous on $(a,b]$, then the graph of $f$ is weakly concave upward on $(a,b]$.
\item $f''\leq 0$ for any $x\in (a,b)$ iff the graph of $f$ is weakly concave downward on $(a,b)$.
\item If $f''\leq 0$ for any $x\in (a,b)$ and $f'$ is continuous on $[a,b)$, then the graph of $f$ is weakly concave downward on $[a,b)$.
\item If $f''\leq 0$ for any $x\in (a,b)$ and $f'$ is continuous on $(a,b]$, then the graph of $f$ is weakly concave downward on $(a,b]$.
\eit
\begin{proof}
By Increasing/Decreasing Test Theorem and First Derivative Test for Concavity Theorem, it is obvious.
\end{proof}
\subsubsection{Inflection Point Theorem}
Let $f\colon I\subseteq\bbR\to\bbR$ be a function. If $(c,f(c))$ is an inflection point of $f$ and $f$ is differentiable at $c$, then $c$ is a critical point of $f'$.
\begin{proof}
By the First Derivative Test for Concavity Theorem, it is obvious.
\end{proof}
\sssc{Extremum of direction derivatives}
For each point $\mb{x}$,
\[\max_{\|\mb{u}\|=1}D_{\mb{u}}\mb{A}(\mb{x})=\|\nabla\mb{A}(\mb{x})\|,\]
\[\arg_{\mb{u}}\max_{\|\mb{u}\|=1}D_{\mb{u}}\mb{A}(\mb{x})=\frac{\nabla\mb{A}(\mb{x})}{\|\nabla\mb{A}(\mb{x})\|},\]
\[\min_{\|\mb{u}\|=1}D_{\mb{u}}\mb{A}(\mb{x})=-\|\nabla\mb{A}(\mb{x})\|,\]
\[\arg_{\mb{u}}\min_{\|\mb{u}\|=1}D_{\mb{u}}\mb{A}(\mb{x})=-\frac{\nabla\mb{A}(\mb{x})}{\|\nabla\mb{A}(\mb{x})\|}.\]
\subsection{Indeterminate forms and L'Hôpital's rule (羅必達法則) or Bernoulli's rule}
\subsubsection{(General form of) L'Hôpital's rule or Bernoulli's rule and indeterminate quotients}
Let $I$ be an open interval and $a\in I$, and $f(x)$ and $g(x)$ be functions differentiable on $I\setminus\{a\}$. If
\bit
\item $\lim_{x\to a}f(x)=\lim_{x\to a}g(x)=0$, where $\lim_{x\to a}\frac{f(x)}{g(x)}$ is called indeterminate form of type $\frac{0}{0}$, or $\abs{\lim_{x\to a}f(x)}=\abs{\lim_{x\to a}g(x)}=\infty$, where $\lim_{x\to a}\frac{f(x)}{g(x)}$ is called indeterminate form of type $\frac{\infty}{\infty}$, and
\item $\lim_{x\to a}\frac{f'(x)}{g'(x)}\in\mathbb{R}\cup\{-\infty,\infty\}$,
\eit
then:
\[\lim_{x\to a}\frac{f(x)}{g(x)}=\lim_{x\to a}\frac{f'(x)}{g'(x)}.\]

Let $f(x)$ and $g(x)$ be functions differentiable on $(a,b)$ with $a<b$. If
\bit
\item $\lim_{x\to a^+}f(x)=\lim_{x\to a^+}g(x)=0$, where $\lim_{x\to a^+}\frac{f(x)}{g(x)}$ is called indeterminate form of type $\frac{0}{0}$, or $\abs{\lim_{x\to a^+}f(x)}=\abs{\lim_{x\to a^+}g(x)}=\infty$, where $\lim_{x\to a^+}\frac{f(x)}{g(x)}$ is called indeterminate form of type $\frac{\infty}{\infty}$, and
\item $\lim_{x\to a^+}\frac{f'(x)}{g'(x)}\in\mathbb{R}\cup\{-\infty,\infty\}$,
\eit
then:
\[\lim_{x\to a^+}\frac{f(x)}{g(x)}=\lim_{x\to a^+}\frac{f'(x)}{g'(x)}.\]

Let $f(x)$ and $g(x)$ be functions differentiable on $(b,a)$ with $a>b$. If
\bit
\item $\lim_{x\to a^-}f(x)=\lim_{x\to a^-}g(x)=0$, where $\lim_{x\to a^-}\frac{f(x)}{g(x)}$ is called indeterminate form of type $\frac{0}{0}$, or $\abs{\lim_{x\to a^-}f(x)}=\abs{\lim_{x\to a^-}g(x)}=\infty$, where $\lim_{x\to a^-}\frac{f(x)}{g(x)}$ is called indeterminate form of type $\frac{\infty}{\infty}$, and
\item $\lim_{x\to a^-}\frac{f'(x)}{g'(x)}\in\mathbb{R}\cup\{-\infty,\infty\}$,
\eit
then:
\[\lim_{x\to a^-}\frac{f(x)}{g(x)}=\lim_{x\to a^-}\frac{f'(x)}{g'(x)}.\]

Let $f(x)$ and $g(x)$ be functions differentiable on $(a,\infty)$. If
\bit
\item $\lim_{x\to\infty}f(x)=\lim_{x\to\infty}g(x)=0$, where $\lim_{x\to\infty}\frac{f(x)}{g(x)}$ is called indeterminate form of type $\frac{0}{0}$, or $\abs{\lim_{x\to\infty}f(x)}=\abs{\lim_{x\to\infty}g(x)}=\infty$, where $\lim_{x\to\infty}\frac{f(x)}{g(x)}$ is called indeterminate form of type $\frac{\infty}{\infty}$, and
\item $\lim_{x\to\infty}\frac{f'(x)}{g'(x)}\in\mathbb{R}\cup\{-\infty,\infty\}$,
\eit
then:
\[\lim_{x\to\infty}\frac{f(x)}{g(x)}=\lim_{x\to\infty}\frac{f'(x)}{g'(x)}.\]

Let $f(x)$ and $g(x)$ be functions differentiable on $(-\infty,a)$. If
\bit
\item $\lim_{x\to -\infty}f(x)=\lim_{x\to -\infty}g(x)=0$, where $\lim_{x\to -\infty}\frac{f(x)}{g(x)}$ is called indeterminate form of type $\frac{0}{0}$, or $\abs{\lim_{x\to -\infty}f(x)}=\abs{\lim_{x\to -\infty}g(x)}=\infty$, where $\lim_{x\to -\infty}\frac{f(x)}{g(x)}$ is called indeterminate form of type $\frac{\infty}{\infty}$, and
\item $\lim_{x\to -\infty}\frac{f'(x)}{g'(x)}\in\mathbb{R}\cup\{-\infty,\infty\}$,
\eit
then:
\[\lim_{x\to -\infty}\frac{f(x)}{g(x)}=\lim_{x\to -\infty}\frac{f'(x)}{g'(x)}.\]
\begin{proof}
Let $f(x)$ and $g(x)$ be functions differentiable on $(a,b)$ with $a<b$. We want to show that if $\lim_{x\to a^+}f(x)=\lim_{x\to a^+}g(x)=0$ and $\lim_{x\to a^+}\frac{f'(x)}{g'(x)}\in\mathbb{R}\cup\{-\infty,\infty\}$, then
\[\lim_{x\to a^+}\frac{f(x)}{g(x)}=\lim_{x\to a^+}\frac{f'(x)}{g'(x)}.\]
Because the existence of the limits, we know that there exists interval $(a,c)$ such that $g(x)$ and $g'(x)$ both $\neq 0$ on $(a,c)$. Let
\[F(x)=\begin{cases}f(x),\quad & x\in (a,c)\\
0,\quad & x=a\end{cases}\]
and
\[G(x)=\begin{cases}g(x),\quad & x\in (a,c)\\
0,\quad & x=a\end{cases}\]
By CMVT, for any $x\in (a,c)$, there exists $c\in (a,x)$ such that
\[\qty(F(x)-F(a))G'(c)=\qty(G(x)-G(a))F'(c).\]
Since $F(a)=G(a)=0$,
\[F(x)G'(c)=G(x)F'(c).\]
Thus,
\[\lim_{x\to a^+}\frac{f(x)}{g(x)}=\lim_{x\to a^+}\frac{F(x)}{G(x)}=\lim_{x\to a^+}\frac{F'(x)}{G'(x)}=\lim_{x\to a^+}\frac{f'(x)}{g'(x)}.\]

Let $f(x)$ and $g(x)$ be functions differentiable on $(a,b)$ with $a<b$. We want to show that if $\lim_{x\to a^+}f(x)=\lim_{x\to a^+}g(x)=\infty$ and $\lim_{x\to a^+}\frac{f'(x)}{g'(x)}\in\mathbb{R}\cup\{-\infty,\infty\}$, then
\[\lim_{x\to a^+}\frac{f(x)}{g(x)}=\lim_{x\to a^+}\frac{f'(x)}{g'(x)}.\]
Because $\lim_{x\to a^+}f(x)=\lim_{x\to a^+}g(x)=\infty$, by quotient law, we know that
\[\lim_{x\to a^+}\frac{1}{f(x)}=\lim_{x\to a^+}\frac{1}{g(x)}=0.\]
Because the existence of the limits, there exists interval $(a,c)$ such that $\frac{1}{f(x)}$ and $\frac{1}{g(x)}$ both $\neq 0$ on $(a,c)$. Let
\[F(x)=\begin{cases}\frac{1}{f(x)},\quad & x\in (a,b)\\
0,\quad & x=a\end{cases}\]
and
\[G(x)=\begin{cases}\frac{1}{g(x)},\quad & x\in (a,b)\\
0,\quad & x=a\end{cases}\]
\[F'(x)=\frac{-f'(x)}{\qty(f(x))^2}.\]
\[G'(x)=\frac{-g'(x)}{\qty(g(x))^2}.\]
By CMVT, for any $x\in (a,c)$, there exists $c\in (a,x)$ such that
\[\qty(F(x)-F(a))G'(c)=\qty(G(x)-G(a))F'(c).\]
Since $F(a)=G(a)=0$,
\[F(x)G'(c)=G(x)F'(c).\]
Thus,
\[\ba
\lim_{x\to a^+}\frac{f(x)}{g(x)}&=\lim_{x\to a^+}\frac{G(x)}{F(x)}=\lim_{x\to a^+}\frac{G'(x)}{F'(x)}\\
&=\lim_{x\to a^+}\frac{g'(x)}{f'(x)}\lim_{x\to a^+}\qty(\frac{f(x)}{g(x)})^2
\ea\]
\[\lim_{x\to a^+}\frac{f(x)}{g(x)}=\lim_{x\to a^+}\frac{f'(x)}{g'(x)}.\]

All other statements can be proven either similarly or by applying other statements.
\end{proof}
\subsubsection{Indeterminate products}
Let $I$ be an open interval and $a\in I$, and $f(x)$ and $g(x)$ be functions differentiable on $I\setminus\{a\}$. If $\lim_{x\to a}f(x)=0$ and $\abs{\lim_{x\to a}g(x)}=\infty$, where $\lim_{x\to a}f(x)g(x)$ is called indeterminate form of type $0\cdot\infty$, we may find $\lim_{x\to a}f(x)g(x)$ by writing the product $f(x)g(x)$ as a quotient:
\[fg=\frac{f}{1/g},\quad\tx{or\ }fg=\frac{g}{1/f}\]
and use general form of L'Hôpital's rule.

Let $f(x)$ and $g(x)$ be functions differentiable on $(a,b)$ with $a<b$. If $\lim_{x\to a^+}f(x)=0$ and $\abs{\lim_{x\to a^+}g(x)}=\infty$, where $\lim_{x\to a^+}f(x)g(x)$ is called indeterminate form of type $0\cdot\infty$, we may find $\lim_{x\to a^+}f(x)g(x)$ by writing the product $f(x)g(x)$ as a quotient:
\[fg=\frac{f}{1/g},\quad\tx{or\ }fg=\frac{g}{1/f}\]
and use general form of L'Hôpital's rule.

Let $f(x)$ and $g(x)$ be functions differentiable on $(b,a)$ with $a>b$. If $\lim_{x\to a^-}f(x)=0$ and $\abs{\lim_{x\to a^-}g(x)}=\infty$, where $\lim_{x\to a^-}f(x)g(x)$ is called indeterminate form of type $0\cdot\infty$, we may find $\lim_{x\to a^-}f(x)g(x)$ by writing the product $f(x)g(x)$ as a quotient:
\[fg=\frac{f}{1/g},\quad\tx{or\ }fg=\frac{g}{1/f}\]
and use general form of L'Hôpital's rule.

Let $f(x)$ and $g(x)$ be functions differentiable on $(a,\infty)$. If $\lim_{x\to\infty}f(x)=0$ and $\abs{\lim_{x\to\infty}g(x)}=\infty$, where $\lim_{x\to\infty}f(x)g(x)$ is called indeterminate form of type $0\cdot\infty$, we may find $\lim_{x\to\infty}f(x)g(x)$ by writing the product $f(x)g(x)$ as a quotient:
\[fg=\frac{f}{1/g},\quad\tx{or\ }fg=\frac{g}{1/f}\]
and use general form of L'Hôpital's rule.

Let $f(x)$ and $g(x)$ be functions differentiable on $(-\infty,a)$. If $\lim_{x\to -\infty}f(x)=0$ and $\abs{\lim_{x\to -\infty}g(x)}=\infty$, where $\lim_{x\to -\infty}f(x)g(x)$ is called indeterminate form of type $0\cdot\infty$, we may find $\lim_{x\to -\infty}f(x)g(x)$ by writing the product $f(x)g(x)$ as a quotient:
\[fg=\frac{f}{1/g},\quad\tx{or\ }fg=\frac{g}{1/f}\]
and use general form of L'Hôpital's rule.
\subsubsection{Indeterminate differences}
Let $I$ be an open interval and $a\in I$, and $f(x)$ and $g(x)$ be functions differentiable on $I\setminus\{a\}$. If $\lim_{x\to a}f(x)=\lim_{x\to a}g(x)\in\{\infty,-\infty\}$, where $\lim_{x\to a}\qty(f(x)-g(x))$ is called indeterminate form of type $\infty-\infty$, we may find $\lim_{x\to a}\qty(f(x)-g(x))$ by writing the difference $f(x)-g(x)$ as a quotient and use general form of L'Hôpital's rule.

Let $f(x)$ and $g(x)$ be functions differentiable on $(a,b)$ with $a<b$. If $\lim_{x\to a^+}f(x)=\lim_{x\to a^+}g(x)\in\{\infty,-\infty\}$, where $\lim_{x\to a^+}\qty(f(x)-g(x))$ is called indeterminate form of type $\infty-\infty$, we may find $\lim_{x\to a^+}\qty(f(x)-g(x))$ by writing the difference $f(x)-g(x)$ as a quotient and use general form of L'Hôpital's rule.

Let $f(x)$ and $g(x)$ be functions differentiable on $(b,a)$ with $a>b$. If $\lim_{x\to a^-}f(x)=\lim_{x\to a^-}g(x)\in\{\infty,-\infty\}$, where $\lim_{x\to a^-}\qty(f(x)-g(x))$ is called indeterminate form of type $\infty-\infty$, we may find $\lim_{x\to a^-}\qty(f(x)-g(x))$ by writing the difference $f(x)-g(x)$ as a quotient and use general form of L'Hôpital's rule.

Let $f(x)$ and $g(x)$ be functions differentiable on $(a,\infty)$. If $\lim_{x\to\infty}f(x)=\lim_{x\to\infty}g(x)\in\{\infty,-\infty\}$, where $\lim_{x\to\infty}\qty(f(x)-g(x))$ is called indeterminate form of type $\infty-\infty$, we may find $\lim_{x\to\infty}\qty(f(x)-g(x))$ by writing the difference $f(x)-g(x)$ as a quotient and use general form of L'Hôpital's rule.

Let $f(x)$ and $g(x)$ be functions differentiable on $(-\infty,a)$. If $\lim_{x\to -\infty}f(x)=\lim_{x\to -\infty}g(x)\in\{\infty,-\infty\}$, where $\lim_{x\to -\infty}\qty(f(x)-g(x))$ is called indeterminate form of type $\infty-\infty$, we may find $\lim_{x\to -\infty}\qty(f(x)-g(x))$ by writing the difference $f(x)-g(x)$ as a quotient and use general form of L'Hôpital's rule.
\subsubsection{Indeterminate powers/exponentials}
Let $I$ be an open interval and $a\in I$, and $f(x)$ and $g(x)$ be functions differentiable on $I\setminus\{a\}$. If $\lim_{x\to a}f(x)=\lim_{x\to a}g(x)=0$, where $\lim_{x\to a}f(x)^{g(x)}$ is called indeterminate form of type $0^0$, we may find $\lim_{x\to a}f(x)^{g(x)}$ by taking natural logarithm:
\[g(x)\ln(f(x)),\]
which is in indeterminate form of $0\cdot\infty$ because $\lim_{x\to a}\ln(f(x))=-\infty$.

Let $I$ be an open interval and $a\in I$, and $f(x)$ and $g(x)$ be functions differentiable on $I\setminus\{a\}$. If $\lim_{x\to a}f(x)=\infty$ and $\lim_{x\to a}g(x)=0$, where $\lim_{x\to a}f(x)^{g(x)}$ is called indeterminate form of type $\infty^0$, we may find $\lim_{x\to a}f(x)^{g(x)}$ by taking natural logarithm:
\[g(x)\ln(f(x)),\]
which is in indeterminate form of $0\cdot\infty$ because $\lim_{x\to a}\ln(f(x))=\infty$.

Let $I$ be an open interval and $a\in I$, and $f(x)$ and $g(x)$ be functions differentiable on $I\setminus\{a\}$. If $\lim_{x\to a}f(x)=1$ and $\lim_{x\to a}g(x)=\infty$, where $\lim_{x\to a}f(x)^{g(x)}$ is called indeterminate form of type $1^{\infty}$, we may find $\lim_{x\to a}f(x)^{g(x)}$ by taking natural logarithm:
\[g(x)\ln(f(x)),\]
which is in indeterminate form of $0\cdot\infty$ because $\lim_{x\to a}\ln(f(x))=0$.

Let $f(x)$ and $g(x)$ be functions differentiable on $(a,b)$ with $a<b$. If $\lim_{x\to a^+}f(x)=\lim_{x\to a^+}g(x)=0$, where $\lim_{x\to a^+}f(x)^{g(x)}$ is called indeterminate form of type $0^0$, we may find $\lim_{x\to a^+}f(x)^{g(x)}$ by taking natural logarithm:
\[g(x)\ln(f(x)),\]
which is in indeterminate form of $0\cdot\infty$ because $\lim_{x\to a^+}\ln(f(x))=-\infty$.

Let $f(x)$ and $g(x)$ be functions differentiable on $(a,b)$ with $a<b$. If $\lim_{x\to a^+}f(x)=\infty$ and $\lim_{x\to a^+}g(x)=0$, where $\lim_{x\to a^+}f(x)^{g(x)}$ is called indeterminate form of type $\infty^0$, we may find $\lim_{x\to a^+}f(x)^{g(x)}$ by taking natural logarithm:
\[g(x)\ln(f(x)),\]
which is in indeterminate form of $0\cdot\infty$ because $\lim_{x\to a^+}\ln(f(x))=\infty$.

Let $f(x)$ and $g(x)$ be functions differentiable on $(a,b)$ with $a<b$. If $\lim_{x\to a^+}f(x)=1$ and $\lim_{x\to a^+}g(x)=\infty$, where $\lim_{x\to a^+}f(x)^{g(x)}$ is called indeterminate form of type $1^{\infty}$, we may find $\lim_{x\to a^+}f(x)^{g(x)}$ by taking natural logarithm:
\[g(x)\ln(f(x)),\]
which is in indeterminate form of $0\cdot\infty$ because $\lim_{x\to a^+}\ln(f(x))=0$.

Let $f(x)$ and $g(x)$ be functions differentiable on $(b,a)$ with $a>b$. If $\lim_{x\to a^-}f(x)=\lim_{x\to a^-}g(x)=0$, where $\lim_{x\to a^-}f(x)^{g(x)}$ is called indeterminate form of type $0^0$, we may find $\lim_{x\to a^-}f(x)^{g(x)}$ by taking natural logarithm:
\[g(x)\ln(f(x)),\]
which is in indeterminate form of $0\cdot\infty$ because $\lim_{x\to a^-}\ln(f(x))=-\infty$.

Let $f(x)$ and $g(x)$ be functions differentiable on $(b,a)$ with $a>b$. If $\lim_{x\to a^-}f(x)=\infty$ and $\lim_{x\to a^-}g(x)=0$, where $\lim_{x\to a^-}f(x)^{g(x)}$ is called indeterminate form of type $\infty^0$, we may find $\lim_{x\to a^-}f(x)^{g(x)}$ by taking natural logarithm:
\[g(x)\ln(f(x)),\]
which is in indeterminate form of $0\cdot\infty$ because $\lim_{x\to a^-}\ln(f(x))=\infty$.

Let $f(x)$ and $g(x)$ be functions differentiable on $(b,a)$ with $a>b$. If $\lim_{x\to a^-}f(x)=1$ and $\lim_{x\to a^-}g(x)=\infty$, where $\lim_{x\to a^-}f(x)^{g(x)}$ is called indeterminate form of type $1^{\infty}$, we may find $\lim_{x\to a^-}f(x)^{g(x)}$ by taking natural logarithm:
\[g(x)\ln(f(x)),\]
which is in indeterminate form of $0\cdot\infty$ because $\lim_{x\to a^-}\ln(f(x))=0$.

Let $f(x)$ and $g(x)$ be functions differentiable on $(a,\infty)$. If $\lim_{x\to\infty}f(x)=\lim_{x\to\infty}g(x)=0$, where $\lim_{x\to\infty}f(x)^{g(x)}$ is called indeterminate form of type $0^0$, we may find $\lim_{x\to\infty}f(x)^{g(x)}$ by taking natural logarithm:
\[g(x)\ln(f(x)),\]
which is in indeterminate form of $0\cdot\infty$ because $\lim_{x\to\infty}\ln(f(x))=-\infty$.

Let $f(x)$ and $g(x)$ be functions differentiable on $(a,\infty)$. If $\lim_{x\to\infty}f(x)=\infty$ and $\lim_{x\to\infty}g(x)=0$, where $\lim_{x\to\infty}f(x)^{g(x)}$ is called indeterminate form of type $\infty^0$, we may find $\lim_{x\to\infty}f(x)^{g(x)}$ by taking natural logarithm:
\[g(x)\ln(f(x)),\]
which is in indeterminate form of $0\cdot\infty$ because $\lim_{x\to\infty}\ln(f(x))=\infty$.

Let $f(x)$ and $g(x)$ be functions differentiable on $(a,\infty)$. If $\lim_{x\to\infty}f(x)=1$ and $\lim_{x\to\infty}g(x)=\infty$, where $\lim_{x\to\infty}f(x)^{g(x)}$ is called indeterminate form of type $1^{\infty}$, we may find $\lim_{x\to\infty}f(x)^{g(x)}$ by taking natural logarithm:
\[g(x)\ln(f(x)),\]
which is in indeterminate form of $0\cdot\infty$ because $\lim_{x\to\infty}\ln(f(x))=0$.

Let $f(x)$ and $g(x)$ be functions differentiable on $(-\infty,a)$. If $\lim_{x\to -\infty}f(x)=\lim_{x\to -\infty}g(x)=0$, where $\lim_{x\to -\infty}f(x)^{g(x)}$ is called indeterminate form of type $0^0$, we may find $\lim_{x\to -\infty}f(x)^{g(x)}$ by taking natural logarithm:
\[g(x)\ln(f(x)),\]
which is in indeterminate form of $0\cdot\infty$ because $\lim_{x\to -\infty}\ln(f(x))=-\infty$.

Let $f(x)$ and $g(x)$ be functions differentiable on $(-\infty,a)$. If $\lim_{x\to -\infty}f(x)=\infty$ and $\lim_{x\to -\infty}g(x)=0$, where $\lim_{x\to -\infty}f(x)^{g(x)}$ is called indeterminate form of type $\infty^0$, we may find $\lim_{x\to -\infty}f(x)^{g(x)}$ by taking natural logarithm:
\[g(x)\ln(f(x)),\]
which is in indeterminate form of $0\cdot\infty$ because $\lim_{x\to -\infty}\ln(f(x))=\infty$.

Let $f(x)$ and $g(x)$ be functions differentiable on $(-\infty,a)$. If $\lim_{x\to -\infty}f(x)=1$ and $\lim_{x\to -\infty}g(x)=\infty$, where $\lim_{x\to -\infty}f(x)^{g(x)}$ is called indeterminate form of type $1^{\infty}$, we may find $\lim_{x\to -\infty}f(x)^{g(x)}$ by taking natural logarithm:
\[g(x)\ln(f(x)),\]
which is in indeterminate form of $0\cdot\infty$ because $\lim_{x\to -\infty}\ln(f(x))=0$.
\subsection{Taylor series}
\subsubsection{Uniqueness}
Let $f\colon D\subseteq\mathbb{R}^m\to\mathbb{R}$ be a function, let $a$ be any point, $I\subseteq D$ be any open ball centered $a$, if there exists a series such that:
\[f(x)=\sum_{\alpha\in\mathbb{N}_0^{\pht{0}m}}\frac{c_{\alpha}}{\alpha !}(x-a)^{\alpha},\]
for all $x\in I$, then
\[c_{\alpha}=D^{\alpha}f(a).\]
\bpr
Fix any multi-index $\beta$. Apply $D^{\beta}$ to both sides
\[D^{\beta}f(x)=\sum_{\alpha\in\bbN_0^{\pht{0}m}}\frac{c_{\alpha}}{\alpha !}D^{\beta}(x-a)^{\alpha}\]
\[D^{\beta}(x-a)^{\alpha}=\bcs
\frac{\alpha !}{(\alpha-\beta)!}(x-a)^{\alpha-\beta},\quad&\beta\leq\alpha\\
0,\quad&\tx{otherwise}\ecs.\]
\[D^{\beta}f(x)=\sum_{\alpha\in\bbN_0^{\pht{0}m}\land\beta\leq\alpha}\frac{c_{\alpha}}{(\alpha-\beta)!}(x-a)^{\alpha-\beta}\]
Evaluate $a$, only $\alpha=\beta$ term survives. Thus,
\[D^{\beta}f(a)=c_{\beta}.\]
\epr
\subsubsection{Taylor's theorem on the reals}
Let $k\in\bbN$, and function $f\colon D\subseteq\bbR\to\bbR$ be $k$ times differentiable at a point $a\in D$. Then there exists a real-valued function $R_k$ defined on an open interval containing $a$ such that
\[f(x)=\sum_{n=0}^k\frac{f^{(n)}(a)}{n!}(x-a)^n+R_k(x),\]
and that
\bit
\item Peano form of remainder:
\[\lim_{x\to a}\frac{R_k(x)}{|x-a|^k}=0.\]
\item Lagrange form of remainder: Suppose $f^{(k+1)}$ exists on $[a,x)$ for $x\geq a$ and $(x,a]$ for $x\leq a$, and $f^{(k)}$ is continuous on $[a,x]$ for $x\geq a$ and $[x,a]$ for $x\leq a$. Then, there exists $\xi$ in $[a,x]$ for $x\geq a$ and $[x,a]$ for $x\leq a$:
\[R_k(x)=\frac{f^{(k+1)}(\xi)}{(k+1)!}(x-a)^{k+1}.\]
\item Integral form of the remainder: Suppose $f^{(k)}$ is absolutely continuous on $[a,x]$ for $x\geq a$ and $[x,a]$ for $x\leq a$. Then
\[R_k(x)=\frac{1}{k!}(x-a)^{k+1}\int_0^1(1-t)^kf^{(k+1)}\qty(a+t(x-a))\dd{t}.\]
\eit
\bpr
PLACEHOLDER
\epr
\subsubsection{Taylor's theorem on higher dimensions}
Let $k,m\in\bbN$, and function $f\colon D\subseteq\bbR^m\to\bbR$ be $k$ times differentiable at a point $a\in D$. Then there exists a real-valued function $R_k$ defined on a neighborhood of $a$ such that
\[f(x)=\sum_{\alpha\in\mathbb{N}_0^{\pht{0}m}\land |\alpha|\leq k}\frac{D^{\alpha}f(a)}{\alpha !}(x-a)^{\alpha}+R_k(x),\]
and that
\bit
\item Peano form of remainder:
\[\lim_{x\to a}\frac{R_k(x)}{|x-a|^k}=0.\]
\item Lagrange form of remainder: Suppose for all $\alpha\in\mathbb{N}_0^{\pht{0}m}\land |\alpha|=k+1$, $f^{\alpha}$ exists on a neighborhood of $a$ whose closure contains $x$, and for all $\alpha\in\mathbb{N}_0^{\pht{0}m}\land |\alpha|=k$, $f^{\alpha}$ is continuous on a neighborhood of $a$ containing $x$. Then, there exists $t\in (0,1)$ such that for $\xi=a+t(x-a)$:
\[R_k(x)=\sum_{\beta\in\mathbb{N}_0^{\pht{0}m}\land |\beta|=k+1}\frac{D^{\beta}f(\xi)}{\beta!}(x-a)^{\beta}.\]
\item Integral form of the remainder: Suppose for all $\alpha\in\mathbb{N}_0^{\pht{0}m}\land |\alpha|=k$, $f^{\alpha}$ is absolutely continuous on a neighborhood of $a$ containing $x$. Then
\[R_k(x)=\sum_{\beta\in\mathbb{N}_0^{\pht{0}m}\land |\beta|=k+1}\frac{k+1}{\beta!}(x-a)^{\beta}\int_0^1(1-t)^k\qty(D^{\beta}f)\qty(a+t(x-a))\dd{t}.\]
\eit
\bpr
PLACEHOLDER
\epr
\subsubsection{Taylor's inequality}
If $\abs{f^{(n+1)}(x)}\leq M$ for $|x-a|\le d$, then the remainder $R_n(x)$ of the Taylor series satisfies the inequality
\[\abs{R_n(x)}\le\frac{M}{(n+1)!}\abs{x-a}^{n+1},\quad\abs{x-a}\le d.\]
\bpr
PLACEHOLDER
\epr
\subsection{Lagrange multiplier (method) (拉格朗日乘數/乘子（法）)}
The Lagrange multiplier method is a method for finding the points where extremes occur of a differentiable function under constraints.
\subsubsection{Univariate form}
Let $f:\,\mathbb{R} \rightarrow \mathbb{R}$ and $g:\,\mathbb{R} \rightarrow \mathbb{R}$ . We want to find the points where extremes of \( f(x) \) under the constraint \( g(x) = \mathbf{0} \) occur. 

First, construct the Lagrangian function \( \mathcal{L}(x,\lambda) \):
\[\mathcal{L}(x,\lambda) = f(x) - \lambda \cdot g(x),\]
where \( \lambda\in\mathbb{R} \) is the Lagrange multiplier (拉格朗日乘數 or 乘子).

Find the derivative of $\mathcal{L}$ and set it to zero:
\[
\dv{\mathcal{L}}{x} = 0
\]
Solve the equation to find \( x \) and \( \lambda \) . 

\text{Statement:} Let $A$ be the set of all solutions for \( x \) such that $\dv{\mathcal{L}}{x} = 0$, and $B$ be the set of all points where extremes of \( f(x) \) under the constraint \( g(x) = \mathbf{0} \) occur. We claim that $B\subseteq A$.
\subsubsection{Multivariate form}
Let \( \mathbf{x} = (x_1, x_2, \dots, x_n) \) be the independent variable vector and $\mathbf{0}$ be the zero vector. Now we have $f:\,\mathbb{R}^n \rightarrow \mathbb{R}$ and $g:\,\mathbb{R}^n \rightarrow \mathbb{R}^c$. We want to find the points where extremes of \( f(\mathbf{x}) \) under the constraint \( g(\mathbf{x}) = \mathbf{0}\) occur. 

First, construct the Lagrangian function \( \mathcal{L}(\mathbf{x},\lambda) \):
\[
\mathcal{L}(\mathbf{x},\lambda) = f(\mathbf{x}) - \lambda \cdot g(\mathbf{x}),
\]
where \( \lambda\in\mathbb{R}^c \) is the Lagrange multiplier vector.

Find the gradient of $\mathcal{L}$ and set it to zero:
\[
\nabla \mathcal{L} = \mathbf{0}
\]
Solve the equation to find \( \mathbf{x} \) and \( \lambda \) . 

\text{Statement:} Let $A$ be the set of all solutions for \( \mathbf{x} \) such that $\nabla \mathcal{L} = \mathbf{0}$, and $B$ be the set of all points where extremes of \( f(\mathbf{x}) \) under the constraint \( g(\mathbf{x}) = \mathbf{0}\) occur. We claim that $B\subseteq A$.
\begin{proof}
Consider $\mathbf{x}^*\in B$. It must satisfy the constraint:
\[g(\mathbf{x}^*) = \mathbf{0}.\]
Any feasible point $\mathbf{x}$ near $\mathbf{x}^*$ must satisfy the constraint. We can represent $\mathbf{x}$ as:
\[\mathbf{x} = \mathbf{x}^* + \delta\mathbf{x},\]
where $\delta\mathbf{x}$ is a differential change tangent to the manifold defined by $g(\mathbf{x})$, that is, it lies in the kernel of $\nabla g(\mathbf{x}^*)$, that is,
\[g(\mathbf{x}^* + \delta\mathbf{x}) = \mathbf{0}.\]
Find the first-order approximation of $f$ at $\mathbf{x}^*$:
\[f(\mathbf{x}^*+ \delta\mathbf{x}) \approx f(\mathbf{x}^*) + \nabla f(\mathbf{x}^*) \cdot \delta\mathbf{x} + O(\|\delta\mathbf{x}\|^2)\]
Find the first-order approximation of $g$ at $\mathbf{x}^*$:
\[g(\mathbf{x}^*+ \delta\mathbf{x}) \approx g(\mathbf{x}^*) + \nabla g(\mathbf{x}^*) \cdot \delta\mathbf{x} + O(\|\delta\mathbf{x}\|^2)\]
Since $\mathbf{x}^*\in B$, for any feasible $\delta\mathbf{x}$ we must have:
\[\nabla f(\mathbf{x}^*) \cdot \delta\mathbf{x} = 0.\]
Since $g(\mathbf{x}^*) = \mathbf{0}$, we have
\[\nabla g(\mathbf{x}^*) \cdot \delta\mathbf{x} = O(\|\delta\mathbf{x}\|^2).\]
Because \( \delta\mathbf{x} \in \ker(\nabla g(\mathbf{x}^)) \), \( \nabla f(\mathbf{x}^) \) can be expressed as a linear combination of \( \nabla g(\mathbf{x}^) \) . This means that there exists a vector to $\lambda$ such that:
\[\nabla\mathcal{L} = \nabla \qty(f(\mathbf{x}) - \lambda \cdot g(\mathbf{x})) = \mathbf{0} \]
\end{proof}
\subsection{Generalized Schwarz's Theorem, Young's Theorem, or Clairaut's Theorem on equality or symmetry of mixed partial derivatives}
For a function $f\colon\Omega\subseteq\mathbb{R}^n\to\mathbb{R}^o$, for any indexed family $A\colon\bbN_{\leq m}\to\bbN_{\leq n}$, for any permutation $\sigma$ of $\bbN_{\leq m}$, define
\[D_{A,\sigma}f=\frac{\partial^mf}{\prod_{k=1}^m\partial x_{A_{\sigma(k)}}.\]
If for all permutations $\sigma$ of $\bbN_{\leq m}$, $D_{A,\sigma}f$ exist on a neighborhood of a fixed point $\mathbf{p}\in\Omega$ and is continuous at $\mathbf{p}$, then all of them are equal at $\mb{p}$.
\subsection{Numerical differentiation (數值微分)}
\subsubsection{Newton's Method (牛頓法) or Newton-Raphson Method (牛頓-拉普森法)}
Newton's method, also known as Newton-Raphson method, is an iterative technique used to approximate the roots of a real-valued function. Given a function $f(x)$ and an initial guess \( x_0 \), Newton's method refines this guess by repeatedly applying the formula:
\[
x_{n+1} = x_n - \frac{f(x_n)}{f'(x_n)},\quad n\in\mathbb{N}
\]
where:
\begin{itemize}
\item \( x_n \) is the current approximation,
\item \( f(x_n) \) is the value of the function at \( x_n \),
\item \( f'(x_n) \) is the derivative of \( f(x) \) evaluated at \( x_n \).
\end{itemize}
PLACEHOLDER: Newton's method converges quadratically or faster if and proof
PLACEHOLDER: Given convergence, suppose that we want the error to be less than a given value $v$, we can take $x_n$ such that $\abs{x_n-x_{n-1}}<v$?